\documentclass[12pt]{article}
\usepackage[margin=1in]{geometry}
\usepackage{setspace}
\usepackage{amsmath,amssymb}
\usepackage{booktabs}
\usepackage{enumitem}
\usepackage{longtable}
\usepackage{array}
\usepackage{xcolor}
\usepackage{hyperref}
\hypersetup{colorlinks=true,linkcolor=blue,urlcolor=blue,citecolor=blue}

\title{Detailed Referee-Style Improvement Report for\\
\textit{Political Polarization and the Interpretation of Auditor Signals}}
\author{Prepared for Neil Hwang}
\date{April 2, 2026}

\begin{document}
\maketitle
\onehalfspacing

\section*{Purpose of This Report}
This report provides a detailed, constructive, and submission-oriented assessment of the current draft titled \textit{Political Polarization and the Interpretation of Auditor Signals}. The goal is to help move the paper from an interesting early-stage idea to a version that is materially more competitive for (i) presentation at the 2026 Conference on Auditing and Capital Markets and (ii) eventual submission to a top accounting journal. The report is written in the style of a developmental editor/referee report rather than a short set of comments.

\section*{Bottom-Line Assessment}
The paper has a \textbf{promising core idea}: investor interpretation of audit-related disclosures may depend on the surrounding political environment, and auditor-change disclosures provide a plausible empirical setting in which identical public signals may generate heterogeneous reactions. This is a good instinct. The topic is novel enough to attract attention, sits at the intersection of auditing, capital markets, and political economy, and is naturally relevant to the PCAOB/CAR conference theme.

At the same time, the current draft is \textbf{not yet ready for submission} to a top accounting journal or to a selective conference in its present form. The paper currently reads as a concept note or an advanced prospectus rather than a finished research paper. The largest issue is not writing polish; it is that the design, identification, measurement, and contribution logic still need substantial strengthening. The draft has a credible \emph{question}, but it does not yet have a sufficiently credible \emph{answering strategy}. In particular, the paper must do more to show that:\
\begin{enumerate}[leftmargin=2em]
    \item the outcome variables capture \emph{interpretation/disagreement} rather than generic news intensity or uncertainty,
    \item the polarization measure is matched to the economically relevant investor environment rather than simply firm headquarters geography,
    \item the empirical design isolates the interpretation channel from omitted firm-level and event-level confounds, and
    \item the paper contributes something distinct to both the auditing literature and the broader finance/political-economy literature.
\end{enumerate}

If these issues are addressed well, the project could become a strong field-archival paper with conference appeal and possibly a credible journal pipeline. Without those changes, reviewers are likely to view the paper as interesting but underidentified and not yet publication-ready.

\section*{Overall Recommendation}
My recommendation is:
\begin{center}
\fbox{\parbox{0.92\textwidth}{
\textbf{Recommendation: Revise aggressively before submission.} Keep the core question, but rebuild the paper around a sharper contribution statement, a more defensible identification strategy, richer event coding, stronger disagreement proxies, and a clearer explanation of why auditor changes are the right setting rather than merely a convenient one.
}}
\end{center}

\section{Why the Paper Has Potential}
The draft already contains several ingredients worth preserving.

\subsection*{1. The research question is interesting}
The idea that political polarization can shape how investors process audit-related signals is timely and potentially important. A large share of prior work in auditing asks whether disclosures are informative \emph{on average}. Your paper asks whether the same disclosure is interpreted differently across environments. That shift in emphasis is potentially valuable.

\subsection*{2. The setting is plausible}
Form 8-K Item 4.01 auditor changes are standardized, salient, and often ambiguous. That ambiguity is useful for a paper about heterogeneous interpretation. The argument that some auditor changes can be benign while others may signal reporting concerns is exactly the sort of ambiguity that could interact with investor priors.

\subsection*{3. The topic fits conference demand}
The project is naturally aligned with auditing, capital markets, and regulation. It is especially relevant for a conference centered on the economic impact of auditing, audit regulation, audit oversight, and capital markets effects. That fit is a real asset.

\subsection*{4. The current draft has a coherent skeleton}
The paper is organized in a recognizable archival-accounting structure: introduction, literature review, theory, hypotheses, design, results placeholders, and conclusion. That means the project already has an architecture that can be upgraded rather than rebuilt from scratch.

\section{Major Issues That Must Be Fixed}
This section identifies the most important problems standing between the current draft and a credible submission.

\subsection{1. The contribution is still too broad and somewhat overstated}
Right now, the introduction makes a large claim: political polarization changes how investors interpret auditor signals, with implications for disclosure regulation and market efficiency. That is potentially true, but the present design is narrower than the claims.

\textbf{Problem.} The paper currently studies one class of events---auditor changes---and proposes to infer a broad mechanism about the interpretation of auditor signals generally. Reviewers may feel the paper is overgeneralizing from a single, noisy event class.

\textbf{What to do.}
\begin{itemize}[leftmargin=2em]
    \item Narrow the headline claim. Do not say you have shown how investors interpret \emph{auditor signals} generally. Say you study \emph{auditor-change disclosures as a setting in which investor interpretation may depend on polarization}.
    \item Make the contribution more precise: the paper is about \emph{cross-sectional variation in market reactions to a regulated audit disclosure}, not about all auditing disclosures.
    \item Reposition the paper as a paper about \emph{belief heterogeneity in audit-related capital markets settings}, using Item 4.01 as the main laboratory.
\end{itemize}

\textbf{Best framing.} A more durable contribution statement would be: \emph{We show that the capital-market consequences of auditor-change disclosures depend systematically on the surrounding political environment, consistent with politically driven differences in interpretation of an ambiguous audit-related signal.}

\subsection{2. The dependent variables do not yet cleanly identify disagreement}
The draft uses $|CAR|$ and abnormal trading volume as the primary outcomes. Those are useful outcomes, but they do not by themselves uniquely identify disagreement.

\textbf{Problem.} A large absolute return can reflect disagreement, but it can also reflect stronger common-direction information. High volume can reflect disagreement, but it can also reflect liquidity demand, arbitrage, short-sale constraints, index-related trading, or event salience. A skeptical reviewer will say: ``You are calling this disagreement, but you are really measuring reaction intensity.''

\textbf{What to do.}
\begin{enumerate}[leftmargin=2em]
    \item Recast $|CAR|$ and abnormal volume as \textbf{primary market-reaction outcomes}, not pure disagreement measures.
    \item Add at least one more direct proxy for disagreement or interpretive dispersion. Candidate measures include:
    \begin{itemize}[leftmargin=2em]
        \item analyst forecast dispersion following the event,
        \item intraday dispersion proxies if available,
        \item bid-ask spread changes as a proxy for differences in valuation/uncertainty,
        \item retail-versus-institutional trading imbalance if feasible,
        \item option-implied disagreement/uncertainty measures if option coverage is sufficient.
    \end{itemize}
    \item Consider a signed-return analysis only after classifying events by expected sign. Otherwise signed CAR tests will look noisy and underpowered.
    \item In the paper, be careful with language: say ``more extreme market reactions and greater trading consistent with more dispersed interpretation'' rather than ``we measure disagreement.''
\end{enumerate}

\subsection{3. Auditor changes are heterogeneous events; your design currently treats them too uniformly}
This is perhaps the most important empirical problem.

\textbf{Problem.} Auditor changes differ drastically in meaning. A resignation with a reported disagreement is not economically comparable to a routine switch, a post-merger change, a move from non-Big 4 to Big 4, or a switch after distress. If these are pooled too loosely, the regression will mix many different information environments.

\textbf{What to do.}
\begin{enumerate}[leftmargin=2em]
    \item Build a serious event-classification scheme from Item 4.01 text. At minimum classify:
    \begin{itemize}[leftmargin=2em]
        \item resignation vs.
        dismissal,
        \item disagreements reported vs.
        none reported,
        \item reportable events disclosed vs.
        none,
        \item Big 4 to non-Big 4, non-Big 4 to Big 4, Big 4 lateral, non-Big 4 lateral,
        \item auditor change around mergers, bankruptcy, or obvious firm restructuring,
        \item switches plausibly attributable to office closure/merger or other mechanical reasons.
    \end{itemize}
    \item Use these event-type variables in both main controls and heterogeneity tests.
    \item Consider restricting the main sample to \textbf{cleaner ambiguous events} where interpretation is actually required, rather than including all 4.01 events.
\end{enumerate}

\textbf{Strong suggestion.} Make event coding a central contribution. A top archival paper in this area needs much more than ``all auditor changes from EDGAR.'' The richness of the coding can become part of the paper's value.

\subsection{4. The key independent variable may be mismatched to the actual investor environment}
The draft measures polarization at the state-year level based on headquarters location. This is simple, but likely too coarse.

\textbf{Problem.} The mechanism is about investor priors and interpretation. But headquarters-state polarization is only an indirect proxy for the beliefs of the marginal investors trading the stock. Reviewers may ask why firm location should determine investor interpretation if ownership is national or global.

\textbf{What to do.}
\begin{enumerate}[leftmargin=2em]
    \item Explicitly justify headquarters geography as capturing the firm's information and stakeholder environment, not literally the ideology of all investors.
    \item Add alternative polarization mappings if possible. Stronger options include:
    \begin{itemize}[leftmargin=2em]
        \item state of major institutional ownership if you can construct it,
        \item local investor base proxies from share turnover/retail participation,
        \item county-level or district-level measures based on headquarters zip code,
        \item media-market level political measures if accessible,
        \item employee or customer political environment proxies for robustness.
    \end{itemize}
    \item Run robustness with finer geography (district or county) and compare to the state measure.
    \item Include a discussion of measurement error. If the state proxy is noisy, say so and frame results as conservative.
\end{enumerate}

\subsection{5. The main regression is not yet conceptually well aligned with the sample}
The current equation includes an \textit{Event} indicator interacted with polarization, but the sample description suggests the dataset may consist only of auditor-change event observations.

\textbf{Problem.} If the sample is only event windows, then the event indicator is redundant. If the sample includes both event and non-event firm-time observations, the design needs much more explanation. As written, the econometric setup is ambiguous.

\textbf{What to do.}
\begin{enumerate}[leftmargin=2em]
    \item Decide whether the paper is:
    \begin{itemize}[leftmargin=2em]
        \item an \textbf{event-only cross-sectional design}, or
        \item a \textbf{panel design with event and non-event windows}.
    \end{itemize}
    \item If event-only, estimate cross-sectional regressions at the event level:
    \[
        Y_{i,e} = \alpha + \beta \, Polarization_{s(i),e} + \Gamma' Controls_{i,e} + FE + \varepsilon_{i,e}.
    \]
    \item If panel, explain exactly how non-event windows are defined, why that comparison is meaningful, and whether event timing is exogenous.
    \item My recommendation is to use an \textbf{event-level design} as the main specification. It is simpler, cleaner, and more credible.
\end{enumerate}

\subsection{6. Identification remains underdeveloped}
At the moment, the draft mostly relies on fixed effects and the idea that polarization varies over time. That is not enough for a top journal.

\textbf{Problem.} Polarization is correlated with many omitted state-level and firm-level factors that can also affect event reactions: media fragmentation, local economic stress, litigation risk, local institutional quality, partisan consumption patterns, industry clustering, and changes in investor composition.

\textbf{What to do.}
\begin{enumerate}[leftmargin=2em]
    \item Add rich event-level controls: prior returns, volatility, size, analyst coverage, institutional ownership, distress, litigation risk, prior restatements, recent earnings news, and contemporaneous 8-K items.
    \item Control for local economic and information environment variables: state GDP growth, unemployment, media fragmentation, election competitiveness, and perhaps social trust proxies if available.
    \item Add industry-year fixed effects or at least event-year and industry fixed effects.
    \item Consider exploiting \textbf{within-firm changes in local political environment around repeated events} when possible.
    \item If feasible, implement a \textbf{placebo test} using non-audit 8-K events of similar salience. If polarization only matters for auditor-change interpretation, that would materially strengthen the mechanism.
    \item Consider an interaction design where ambiguity is pre-classified. The mechanism predicts stronger effects where interpretation is required, not uniformly across all events.
\end{enumerate}

\subsection{7. The theory section is directionally useful but not yet publication-level}
The theory is serviceable as intuition, but it will not yet satisfy a demanding referee.

\textbf{Problem.} The model currently assumes the conclusion: polarization widens priors, which widens posterior dispersion, which increases trade. That logic is fine, but the presentation is thin. The proposition proof is too informal, and the theory does not sufficiently connect to the observable empirical proxies.

\textbf{What to do.}
\begin{enumerate}[leftmargin=2em]
    \item Reframe the theory section as a parsimonious organizing framework, not a formal theoretical contribution unless you are willing to deepen it substantially.
    \item Tighten the proposition language. Do not state a general mean-preserving spread result without care unless you can prove it cleanly or cite the appropriate result.
    \item Link each theoretical object to an empirical proxy: posterior dispersion $\rightarrow$ disagreement/reactivity measures; lower signal precision $\rightarrow$ no disclosed reason, no disagreement detail, non-Big 4 switches, complex firms, etc.
    \item Clarify why polarization changes priors about auditor credibility or institutional trust specifically, not just generic risk appetite.
\end{enumerate}

\subsection{8. The literature review is too thin for a top journal submission}
The current literature review is competent but abbreviated.

\textbf{Problem.} Right now the paper cites broad areas, but it does not yet convince the reader that you fully understand where this paper fits and how it is distinct. Top-journal reviewers will expect a much more developed discussion of (i) auditor-change literature, (ii) market reactions to audit signals, (iii) belief heterogeneity and disagreement, (iv) political ideology in markets, and (v) institutional trust and disclosure interpretation.

\textbf{What to do.}
\begin{enumerate}[leftmargin=2em]
    \item Expand the literature review substantially.
    \item Separate \textbf{contribution} from \textbf{relatedness}. Many papers are related; only a few are true contribution margins.
    \item Add a short subsection specifically on \emph{interpretation of ambiguous disclosures}. That is where your project really lives.
    \item Add a short subsection on why auditing is a special institutional setting: standardized disclosure, delegated assurance, regulatory infrastructure, and trust in gatekeepers.
\end{enumerate}

\subsection{9. The introduction currently reads as if results already exist, but the draft is not yet evidence-complete}
The introduction includes placeholders for results and economic magnitudes.

\textbf{Problem.} This creates a mismatch in tone. The paper sounds finished, but the evidence sections are skeletal. For conference review and journal review, this mismatch will reduce credibility.

\textbf{What to do.}
\begin{itemize}[leftmargin=2em]
    \item Do not write as though the paper is complete until the tables exist.
    \item Either complete the empirical work and then write a full introduction, or temporarily adopt a more measured tone.
    \item Once estimates are in hand, rewrite the introduction around exact economic magnitudes, sample size, and your two strongest identification/robustness points.
\end{itemize}

\section{Section-by-Section Comments}

\subsection{Title}
Current title: \textit{Political Polarization and the Interpretation of Auditor Signals}

\textbf{Assessment.} Strong topic signal, but slightly broader than the actual design.

\textbf{Recommendation.} Consider one of the following:
\begin{itemize}[leftmargin=2em]
    \item \textit{Political Polarization and Market Reactions to Auditor-Change Disclosures}
    \item \textit{Political Polarization and the Interpretation of Auditor-Change Signals}
    \item \textit{Auditor Changes, Political Polarization, and Investor Interpretation}
\end{itemize}
These titles still sound ambitious, but they better match the setting.

\subsection{Abstract}
\textbf{What works.} The abstract clearly identifies the question, mechanism, setting, and main intended outcomes.

\textbf{What needs improvement.}
\begin{itemize}[leftmargin=2em]
    \item The abstract currently promises specific results that are not yet shown in the draft.
    \item It overstates certainty relative to the stage of the evidence.
    \item It needs concrete numbers once results are available.
\end{itemize}

\textbf{Recommendation.} Rewrite the abstract last. A top abstract should include: question, setting, sample period, main measure, key coefficient magnitude, one mechanism/heterogeneity result, and one sentence on contribution.

\subsection{Introduction}
\textbf{What works.}
\begin{itemize}[leftmargin=2em]
    \item The opening intuition is understandable.
    \item The choice of auditor changes is motivated reasonably well.
    \item The paper identifies a plausible mechanism tied to heterogeneous priors.
\end{itemize}

\textbf{What needs improvement.}
\begin{itemize}[leftmargin=2em]
    \item The introduction is still too general and repeats broad claims.
    \item It needs a sharper ``why should we care'' paragraph for the auditing audience.
    \item It should explain why auditor-change events are especially suited to reveal interpretive heterogeneity.
    \item It should explain why the effect is not simply stronger reactions in more uncertain or distressed states.
\end{itemize}

\textbf{Recommendation.} Rewrite the introduction around four paragraphs:
\begin{enumerate}[leftmargin=2em]
    \item Big question: standardized audit disclosures may not be interpreted uniformly.
    \item Why auditor changes are a strong setting: mandatory, salient, often ambiguous, and institution-dependent.
    \item What you do: code Item 4.01 events, measure polarization, test market reactions and heterogeneity.
    \item What you find and why it matters: one sentence each on main result, ambiguity, and contribution.
\end{enumerate}

\subsection{Related Literature}
\textbf{Assessment.} This section is a placeholder-level literature review, not yet a journal-quality one.

\textbf{Recommendation.} Expand into a more strategic map:
\begin{enumerate}[leftmargin=2em]
    \item Auditor changes and market reactions.
    \item Audit disclosure informativeness and capital-market consequences.
    \item Political beliefs, trust, and financial-market interpretation.
    \item Heterogeneous beliefs, disagreement, and trading.
    \item Your contribution paragraph: what exactly no prior paper has done.
\end{enumerate}

\subsection{Theory and Hypotheses}
\textbf{Assessment.} The current framework is helpful as intuition but should either be strengthened or simplified.

\textbf{Recommendation.}
\begin{itemize}[leftmargin=2em]
    \item Keep a compact conceptual model unless you intend to produce a deeper theoretical section.
    \item Convert propositions into clean, testable implications.
    \item Merge overlapping hypotheses. Right now H1 and H3 are fine, but H2 may be difficult unless event sign is credibly classified.
\end{itemize}

A better hypothesis set may be:
\begin{description}[leftmargin=2em,style=nextline]
    \item[H1:] Market reactions to auditor-change disclosures are more extreme in more polarized environments.
    \item[H2:] The relation in H1 is stronger when the disclosure is more ambiguous.
    \item[H3:] The relation in H1 is stronger when the event is more consequential for perceived reporting credibility.
\end{description}

\subsection{Empirical Design}
\textbf{Assessment.} This is the section that most needs redevelopment.

\textbf{Specific improvements needed.}
\begin{enumerate}[leftmargin=2em]
    \item Clarify sample construction in detail.
    \item Explain the exact event date used (filing date, acceptance timestamp, next trading day if after hours, etc.).
    \item Justify the event window and estimation window.
    \item Address confounding simultaneous disclosures in the same 8-K or nearby filings.
    \item Add event-type coding and ambiguity measures.
    \item Re-specify the regression at the event level.
\end{enumerate}

\subsection{Results}
\textbf{Assessment.} The results section currently contains placeholders only.

\textbf{Recommendation.} The final results section should be organized as follows:
\begin{enumerate}[leftmargin=2em]
    \item Descriptive statistics and event taxonomy.
    \item Baseline event-study evidence.
    \item Main cross-sectional regressions.
    \item Ambiguity tests.
    \item Event-type heterogeneity.
    \item Mechanism/disagreement proxies.
    \item Robustness and placebo tests.
\end{enumerate}

\section{Most Likely Referee Critiques and How to Preempt Them}
Below are the criticisms you are most likely to face, together with how to answer them in the paper.

\subsection*{Referee concern 1: ``This is just a generic uncertainty paper.''}
\textbf{Response strategy:} Show that the effect is strongest in ex ante ambiguous auditor-change events, not uniformly across all news. Add placebo tests using other 8-K items or non-audit disclosures.

\subsection*{Referee concern 2: ``State-level polarization at headquarters is not investor polarization.''}
\textbf{Response strategy:} Add finer geography and alternative mappings; emphasize that headquarters polarization proxies the local information and institutional environment in which the disclosure is interpreted and discussed.

\subsection*{Referee concern 3: ``$|CAR|$ and volume are not disagreement.''}
\textbf{Response strategy:} Soften the claim and add at least one more direct disagreement proxy.

\subsection*{Referee concern 4: ``Auditor changes differ too much to pool.''}
\textbf{Response strategy:} Build a rich event taxonomy and show results within cleaner subsets.

\subsection*{Referee concern 5: ``The mechanism is political polarization, but maybe this is just local economic distress or institutional weakness.''}
\textbf{Response strategy:} Add controls for local economic conditions, institutional quality, media fragmentation, and firm distress. Then show the result survives.

\subsection*{Referee concern 6: ``The paper does not show why this matters for auditing specifically.''}
\textbf{Response strategy:} Emphasize delegated assurance, standardized regulatory disclosure, investor trust in audit oversight, and the role of ambiguity in auditor-change filings.

\section{Concrete Empirical Upgrades Needed Before Submission}

\subsection{A. Build a much richer event file}
For each auditor-change event, manually or algorithmically code:
\begin{itemize}[leftmargin=2em]
    \item event date and filing timestamp,
    \item resignation/dismissal,
    \item stated reason,
    \item disagreement indicator,
    \item reportable event indicator,
    \item predecessor and successor auditor names,
    \item Big 4 status change,
    \item whether the firm disclosed consultation with the new auditor,
    \item whether the firm had recent restatement, going concern, internal control issues, or unusual accruals,
    \item whether the event coincides with M\&A, bankruptcy, or major financing.
\end{itemize}

\subsection{B. Add better measures of ambiguity}
Your current paper says the mechanism is stronger for ambiguous signals. That needs a serious operationalization. Possible ambiguity indicators:
\begin{itemize}[leftmargin=2em]
    \item no disclosed reason for change,
    \item no disagreement but abrupt resignation,
    \item textual vagueness in Item 4.01 language,
    \item no change in Big 4 status,
    \item complex or opaque firms (high accruals, high intangibles, low analyst coverage),
    \item prior audit issues that make interpretation less clear.
\end{itemize}

\subsection{C. Improve the outcome side}
At minimum, estimate:
\begin{itemize}[leftmargin=2em]
    \item $CAR[-1,+1]$,
    \item $|CAR[-1,+1]|$,
    \item abnormal turnover/volume,
    \item change in bid-ask spread,
    \item post-event analyst dispersion if data are available,
    \item longer-window drift if the interpretation effect unfolds gradually.
\end{itemize}

\subsection{D. Add strong robustness tests}
Recommended robustness analyses:
\begin{itemize}[leftmargin=2em]
    \item alternative event windows,
    \item alternative polarization measures,
    \item state versus district versus county geographic mapping,
    \item excluding single-district states,
    \item excluding crisis years,
    \item excluding firms with contemporaneous major disclosures,
    \item clustering by firm and state or two-way clustering where appropriate,
    \item winsorization and median regressions,
    \item industry-year fixed effects.
\end{itemize}

\subsection{E. Add placebo and falsification tests}
These could be especially valuable:
\begin{itemize}[leftmargin=2em]
    \item placebo on non-audit 8-K events,
    \item placebo on clearly unambiguous auditor changes,
    \item placebo using future polarization,
    \item placebo outcomes outside the event window.
\end{itemize}

\section{How to Improve the Paper for the PCAOB/CAR Conference Specifically}
A conference paper is not identical to a journal paper. For this conference, you should emphasize not only novelty but also why the project matters for auditing policy and capital markets.

\subsection*{Conference-facing strengths to emphasize}
\begin{itemize}[leftmargin=2em]
    \item The project studies a \textbf{regulated audit disclosure} rather than a generic political-finance outcome.
    \item The paper speaks directly to \textbf{capital-market consequences of auditing information}.
    \item The mechanism connects to \textbf{how investors process institution-backed assurance signals}.
    \item The topic is relevant to debates about whether additional disclosure always reduces uncertainty.
\end{itemize}

\subsection*{What conference reviewers will want to see}
\begin{itemize}[leftmargin=2em]
    \item completed baseline tables,
    \item a convincing event taxonomy,
    \item one or two strong mechanism tests,
    \item clear economic magnitudes,
    \item a concise statement of policy relevance without overclaiming.
\end{itemize}

\subsection*{How to pitch the paper in the cover email or abstract}
The best pitch is not ``politics matters for everything.'' The best pitch is:
\begin{quote}
We examine whether political polarization changes how investors interpret auditor-change disclosures. Using Item 4.01 filings, we test whether identical audit-related public signals produce more extreme market reactions in more polarized environments, particularly when the disclosure is more ambiguous. The findings speak to the capital-market consequences of audit communication and the limits of disclosure-based regulation when investor priors differ systematically.
\end{quote}

\section{How to Improve the Paper for a Top Accounting Journal}
For a top journal, the bar is higher. You will need more than an interesting conference result.

\subsection*{What top-journal reviewers will likely require}
\begin{enumerate}[leftmargin=2em]
    \item A clearly identified mechanism and not just a suggestive pattern.
    \item A much more credible design linking polarization to interpretation rather than to confounds.
    \item A richer treatment of alternative explanations.
    \item Stronger placement in the auditing literature.
    \item Writing that is tighter, less repetitive, and less promotional.
\end{enumerate}

\subsection*{The most important strategic decision}
You need to decide whether the paper's natural home is:
\begin{itemize}[leftmargin=2em]
    \item an \textbf{auditing journal} with strong capital-markets orientation, or
    \item a broader \textbf{top accounting journal} where the contribution must travel beyond audit specialists.
\end{itemize}

If you target a very top generalist accounting outlet, the paper must make clear why the result changes how we think about disclosure interpretation or assurance value more broadly. If you target a strong auditing-oriented outlet first, the paper can be more tightly focused on auditor-change disclosures and regulatory implications.

\section{Writing Improvements Needed}
The writing is already reasonably clear, but several improvements are needed.

\subsection*{1. Reduce overclaiming}
Avoid sentences that suggest the paper has already established broad policy conclusions unless the evidence directly supports them.

\subsection*{2. Be more precise}
Replace vague statements like ``in economically meaningful ways'' with actual effect sizes once available.

\subsection*{3. Use tighter transitions}
Several paragraphs currently restate the same intuition in slightly different ways. Compress repetition and use the saved space for identification and measurement detail.

\subsection*{4. Distinguish carefully among related concepts}
Do not conflate:
\begin{itemize}[leftmargin=2em]
    \item disagreement,
    \item uncertainty,
    \item ambiguity,
    \item salience,
    \item information asymmetry,
    \item reaction intensity.
\end{itemize}
A strong paper will define these distinctly and explain which one is observed versus inferred.

\section{Recommended Revised Paper Structure}
Here is a stronger structure for the next full draft.

\begin{enumerate}[leftmargin=2em]
    \item \textbf{Introduction} --- sharpened contribution, exact results, exact magnitude.
    \item \textbf{Institutional Background and Setting} --- Item 4.01 disclosure requirements; why auditor changes are informative yet ambiguous.
    \item \textbf{Conceptual Framework and Hypotheses} --- concise model plus clean hypotheses.
    \item \textbf{Sample, Event Coding, and Variable Construction} --- this should be much richer than the current design section.
    \item \textbf{Baseline Event-Study Evidence}.
    \item \textbf{Main Cross-Sectional Tests}.
    \item \textbf{Ambiguity and Mechanism Tests}.
    \item \textbf{Alternative Explanations and Robustness}.
    \item \textbf{Conclusion and Policy Implications}.
\end{enumerate}

\section{Priority-Ordered Revision Plan}
If you want the highest return on effort, revise in this order.

\subsection*{Tier 1: Essential before any submission}
\begin{enumerate}[leftmargin=2em]
    \item Finalize the sample and event coding.
    \item Redesign the main empirical specification.
    \item Produce real baseline tables and figures.
    \item Add ambiguity and event-type heterogeneity tests.
    \item Rewrite the introduction once the evidence exists.
\end{enumerate}

\subsection*{Tier 2: Essential for stronger conference/journal credibility}
\begin{enumerate}[leftmargin=2em]
    \item Add direct disagreement or uncertainty proxies beyond $|CAR|$ and volume.
    \item Add alternative polarization mappings and richer controls.
    \item Add placebo/falsification tests.
    \item Expand the literature review materially.
\end{enumerate}

\subsection*{Tier 3: Important for a top-journal push}
\begin{enumerate}[leftmargin=2em]
    \item Improve the theory section or simplify it strategically.
    \item Add stronger economic-magnitude interpretation.
    \item Refine policy implications so they are disciplined and not overstated.
\end{enumerate}

\section{Illustrative Tables and Figures You Should Add}
\subsection*{Tables}
\begin{itemize}[leftmargin=2em]
    \item Table 1: Sample selection and event taxonomy.
    \item Table 2: Summary statistics.
    \item Table 3: Event-study returns and volume by event type.
    \item Table 4: Main regression results.
    \item Table 5: Ambiguity tests.
    \item Table 6: Alternative measures of polarization and geography.
    \item Table 7: Alternative explanations and placebo tests.
\end{itemize}

\subsection*{Figures}
\begin{itemize}[leftmargin=2em]
    \item Figure 1: Timeline of sample construction.
    \item Figure 2: Distribution of event types.
    \item Figure 3: Map or distribution of polarization measure.
    \item Figure 4: Binned scatter of event reaction versus polarization.
    \item Figure 5: Heterogeneity plot for ambiguity classes.
\end{itemize}

\section{Editorial Verdict}
As a developmental referee, my view is that this project is \textbf{worth pursuing}. The central idea is good, the setting is promising, and the conference fit is real. But the current draft is still at a stage where the novelty of the question exceeds the credibility of the evidence. The revision task is therefore not cosmetic. It is architectural.

If you make the changes above, especially on event coding, identification, disagreement measurement, and contribution framing, the project can become much stronger. Without those changes, the paper will likely be seen as a well-written but still preliminary idea paper.

\section*{Concise Action Summary}
\begin{center}
\fbox{\parbox{0.94\textwidth}{
\textbf{Keep:} the core question, the audit-setting intuition, and the polarization/disagreement mechanism.\\[0.4em]
\textbf{Fix urgently:} event heterogeneity, outcome interpretation, geography mapping, identification, and literature positioning.\\[0.4em]
\textbf{Best next step:} rebuild the empirical section around a richly coded event-level dataset and then rewrite the introduction and contribution statement around the evidence that survives the strongest tests.
}}
\end{center}

\end{document}
